% Detailed comparison tables retained for implementers and auditors.
% The main chapters use compact summaries first and explain how to read them.
\chapter{Detailed comparison tables}
\label{app:detailed-tables}

The tables in this appendix preserve the complete comparison records behind the
shorter tables in Chapters~\ref{ch:languages}, \ref{ch:ensemble} and
\ref{ch:evaluation}. They are lookup material. A reader who wants the argument
can remain in the main chapters; an implementer can use these tables to recover
the full row-by-row specification.

\section{Ensemble responsibilities}
\begin{table}[htbp]
\centering\small
\caption{Responsibilities that create different opportunities to detect error.}
\label{tab:ensemble-roles-detail}
\begin{tabularx}{\textwidth}{P{31mm}YY}
\toprule
Role & Question answered & Failure that must remain visible\\
\midrule
Source inventory & What material is present and incorporated? & Missing attachment, unmatched paragraph, inaccessible dependency\\
Normative reader & What obligations, conditions and exceptions are expressed? & Omitted modality, scope or qualification\\
Alternative reader & What supported choice changes a consequence? & No supported alternative found; unexplored family\\
Argument critic & What defeats the premises or inference? & Unanswered objection or unverified authority\\
Formal checker & What follows from this specified rule? & Invalid encoding, unknown solver result, limited domain\\
Reference reviewer & What is expected for this scenario, and why? & Disagreement, incomplete facts, shared authorship\\
Java checker & Does this package preserve the accepted specification? & Mismatch, unsupported operation, stale build\\
Dossier assembler & Are all required results and dispositions present? & Missing member, open issue, broken evidence link\\
\bottomrule
\end{tabularx}
\end{table}

\section{Resolution evidence}
\begin{longtable}{@{}P{27mm}P{56mm}P{65mm}@{}}
\caption{Discrepancy-specific actions and closure evidence.}\label{tab:resolution-matrix-detail}\\
\toprule
Discrepancy & Appropriate next action & What can justify closure\\\midrule\endfirsthead
\toprule Discrepancy & Appropriate next action & What can justify closure\\\midrule\endhead
Missing source unit & Compare original rendering and independent inventory & Retained unit, corrected inventory and reviewed disposition; agreement among summaries is insufficient.\\
Unavailable dependency & Retrieve the named version or seek authorized clarification & Inspected applicable source, or an explicit adjudication of the remaining question; a guessed definition is insufficient.\\
Different source versions & Establish the assessment cut-off and amendment relationship & A documented version selection and reconsideration of affected candidates.\\
Actor or activity scope & Inspect addressee, operative language and profile facts & Reviewed applicability proposition with evidence for the actual bank activity.\\
% BEGIN REVISION ADDITION gap-resolution-table
\bottomrule
\end{longtable}
\begingroup\normalsize
\begin{figure}[H]
\centering
\TeachingFlow{Classify the discrepancy}{Obtain the relevant evidence}{Close only the question answered}
\caption{Resolution requires evidence of the appropriate kind.}\label{fig:gap-resolution-table}
\end{figure}

\endgroup
\begin{longtable}{@{}P{27mm}P{56mm}P{65mm}@{}}

\toprule Discrepancy & Appropriate next action & What can justify closure\\\midrule\endhead
% END REVISION ADDITION gap-resolution-table
Disputed classification & Identify definition and distinguishing factual evidence & An accepted definition and sufficient facts, or an attributable adjudication; a Boolean label alone is insufficient.\\
Exception attachment & Construct component and whole-arrangement examples & Source-grounded structural choice, with contrary reading disposition and regression cases.\\
Necessary versus sufficient condition & Compare implication directions on boundary scenarios & A reviewed direction of dependence and a rule that preserves other required conditions.\\
Timing or transition & Separate publication, commencement, valid time and knowledge time & Reviewed interval and transition rule with endpoint examples.\\
Factual conflict & Obtain correction or apply an approved evidence-priority rule & Attributable supersession or reviewed normalization; newest arrival alone is insufficient.\\
Formal mismatch & Produce and replay a feasible difference witness & Validated repair or domain-qualified equivalence, with source assumptions retained.\\
Unsupported operation & Inspect whether a supported encoding preserves meaning & Reviewed encoding plus conformance evidence, or an approved runtime extension; compiler rejection does not defeat the reading.\\
Java mismatch & Compare full semantic results and trace & Corrected exact package with repeated affected verification and new build identity.\\
Missing ensemble member & Retry under budget or use an authorized replacement & Required role output with documented independence limits; silence is not agreement.\\
Search truncation & Investigate retained frontier or narrow the claimed scope & Completed declared finite search or explicit reviewed restriction; a high best score is insufficient.\\
Broken approval lineage & Reconstruct exact source, report, bundle and build links & Fresh valid approvals for exact identities; a matching display title is insufficient.\\
\bottomrule
\end{longtable}


\section{Software and standards considered}
\begingroup\small
\begin{longtable}{P{0.19\textwidth}P{0.37\textwidth}P{0.34\textwidth}}
\caption{Candidate software and standards, with a deliberately bounded role.}
\label{tab:software-detail}\\
\toprule Candidate & Contribution to LegalMath & Limit or prototype decision \\\midrule
\endfirsthead
\toprule Candidate & Contribution to LegalMath & Limit or prototype decision \\\midrule
\endhead
Catala & Literate definitions, exceptions, executable decisions and verification-condition ideas. & Compare one adapter on the SPI decision fragment; do not assume event-workflow coverage or end-to-end verified compilation. \citep{catalacode}\\
DMN 1.5; Apache KIE / Drools & Standard decision tables and Java-oriented decision execution. & Preserve hit policies, explicit unknown handling and explanations. Add a separate event model. \citep{dmn2024,droolscode}\\
LegalRuleML & Source, authority, interpretation and normative metadata for interchange. & Adopt a small profile; a reasoner and execution semantics are separate components. \citep{legalruleml2021}\\
Stipula; Stipula--KeY & Contract states, permissions, scheduled events and a Java/JML proof experiment. & Research adapter for a small workflow; document event abstraction and loop restrictions. \citep{stipulacode,stipulakeycode}\\
eFLINT & Normative simulation, duties and violations, including actions that occur despite a prohibition. & Borrow concepts and compare a bounded scenario; pin the implementation and semantics version. \citep{eflint2026}\\
L4 & Legal DSL, regulative rules, event traces, IDE and generated decision services. & Candidate for a later authoring/core comparison; documentation does not establish equivalence with RuleIR. \citep{l4code}\\
Symboleo & Obligation/power lifecycles and a model-checking toolchain. & State-model alternative; preliminary paper case and software inspected, later journal paper remains a retrieval gap. \citep{symboleo2020,symboleocode}\\
Regorous / PCL & Defeasible norms connected to annotated business processes. & Pilot evidence supports the method; deployed version, availability and license remain adoption questions. \citep{semanticcompliance2016}\\
Blawx / s(CASP) & Visual rule authoring, explanations and hypothetical reasoning. & Authoring experiment; Blawx describes its current scope as educational and experimental. \citep{blawxcode,slaw2024}\\
Z3; existing Lean route & Counterexamples, satisfiability, selected invariants and checked proof targets. & Specify supported theories and assumptions; preserve timeout and unknown outcomes. \citep{z3code,localrepos}\\
OpenFisca & Versioned computable rules and a tax/benefit microsimulation architecture. & Useful design analogue; the primary problem here includes conduct and workflow duties. \citep{openfiscacode}\\
Open Policy Agent & Rego policies evaluated against structured data; potential deployment adapter. & Its undefined-value behavior must be reconciled with the legal rule semantics. It supplies no natural-language legal interpretation. \citep{opadocs}\\
Accord Project / Cicero & Adjacent work on digital contracts and reusable contract technology. & Monitor for document/template integration; a fixed contract template differs from an amended regulatory corpus. \citep{cicerocode}\\
\bottomrule
\end{longtable}
\endgroup


\section{Cross-circular acceptance cases}
\begingroup\small
\begin{longtable}{P{0.09\textwidth}P{0.41\textwidth}P{0.4\textwidth}}
\caption{Initial acceptance scenarios for the selected scope.}
\label{tab:cases-detail}\\
\toprule Case & Synthetic fact or source change & Required demonstration \\\midrule
\endfirsthead
\toprule Case & Synthetic fact or source change & Required demonstration \\\midrule
\endhead
F1 & Portfolio HK\$40m; qualifying net assets HK\$0. & Establish the financial
condition via the inclusive portfolio comparison.\\
F2 & Portfolio HK\$35m; qualifying net assets HK\$90m. & Establish the financial
condition via the alternative net-asset route.\\
F3 & Portfolio unknown; qualifying net assets HK\$70m. & Leave the financial
condition unresolved; identify missing portfolio evidence.\\
F4 & Portfolio unknown; qualifying net assets HK\$80m. & Establish the financial
condition via the known sufficient route.\\
F5 & Assets HK\$120m, liabilities HK\$15m, primary residence HK\$30m. & Derive
HK\$75m for the illustrated net-asset calculation; assess the separate portfolio
route independently.\\
O1 & HK\$60m joint portfolio with a non-associate; written client allocation 20\%. &
Use HK\$12m for this portfolio-share illustration; explain the agreement and actor roles.\\
K1 & Transaction-experience route established only for one category; a different
category proposed. & Do not propagate a global experience flag; inspect the
alternative qualification routes and selected categories.\\
C1 & A qualifying wealth record, but conservative investment objectives. & Explain
the investment-objective exclusion; wealth alone does not establish SPI treatment.\\
E1 & Designated-account monitoring selected; exposure rises above threshold without
new funds. & Show the relevant FAQ conditions and restrictions; avoid a universal
transaction-block rule.\\
E2 & A withdrawal instruction precedes a proposed streamlined decision in the
adopted event order. & Use the changed consent state and retain the original
instruction and ordering evidence.\\
T1 & A third-party verification clause with an SFC-request trigger. & Preserve the
trigger; distinguish an absent request from missing request records.\\
T2 & Old tokenisation source replaced by the 2026 circular. & Preserve both editions,
record replacement, identify affected rules and reproduce historical decisions.\\
M1 & Product-linked gift versus a discount of fees or charges. & Preserve the stated
exception and the scope of the cited marketing provision.\\
M2 & Daily dealing fund with weekly redemption versus a different daily cut-off. &
Distinguish reduced dealing frequency from administrative procedure.\\
D1 & A thematic-review announcement enters the source collection. & Classify and
track the announcement without inventing a new numerical eligibility criterion.\\
\bottomrule
\end{longtable}
\endgroup


\section{Evaluation evidence contract}
\begingroup\small
\begin{longtable}{P{0.25\textwidth}P{0.39\textwidth}P{0.26\textwidth}}
\caption{Proposed evidence contract for the prototype evaluation. These are
prospective criteria, not results.}\label{tab:evaluation-detail}\\
\toprule Evidence & Definition and use & Decision role \\\midrule
\endfirsthead
\toprule Evidence & Definition and use & Decision role \\\midrule
\endhead
Material interpretation and decision defects & Independently adjudicated errors
in source coverage, rule meaning, resulting decisions or obligations. No unresolved
critical defect may be carried into a proposed operational scope. & Primary quality
criterion; blocks promotion of the affected scope.\\
Review effort & End-to-end person-time for a completed, accepted change, including
source work, corrections, escalations and reusable-rule preparation. & Primary
business criterion, conditional on quality.\\
Unknown and review outcomes & Frequency, reason, downstream handling and time to
resolution; distinguish necessary abstention from unhelpful failure. & Explanatory
diagnostic and business cost; incorrect silent resolution is a quality defect.\\
Source and dependency completeness & Every in-scope rule has inspected support and
resolved imports; every required selected clause is accounted for. & Release veto
for the selected corpus. Does not establish completeness of all HK regulation.\\
Adapter agreement and mutations & Supported semantics agree; targeted changes in
operators, scope, dates and exceptions are detected by the required checks. &
Engineering veto on the adapter; mutation scores alone do not establish legal fidelity.\\
Reproducibility & A saved rule version, fact snapshot and event sequence reproduce
the original result and explanation references. & Engineering veto on historical replay.\\
Latency and token/runtime cost & Measure on the specified local workload, including
failed drafts and review retries. & Explanatory and operating-budget evidence; no
quality trade-off is implied.\\
\bottomrule
\end{longtable}
\endgroup


\section{Operational readiness evidence}
\begin{longtable}{@{}P{31mm}P{60mm}P{57mm}@{}}
\caption{Evidence required for a bounded deployment decision.}\\
\toprule
Area & Evidence to inspect & Question that remains a human responsibility\\\midrule\endfirsthead
\toprule Area & Evidence to inspect & Question that remains a human responsibility\\\midrule\endhead
Meaning & Source inventory, candidates, objections, adjudication & Is this reading justified for the selected activity and time?\\
Data & Fact dictionary, provenance, validity, conflict policy & Do the bank's records establish the facts the rule needs?\\
Execution & Conformance, witnesses, mutations, Java comparison & Is the supported implementation scope adequate for the process?\\
Authority & Exact manifests, reviewer roles, signatures & Are the required people authorized and accountable for this release?\\
Operation & Host dispositions, race tests, fallback and incident rehearsal & Can staff act correctly when the result is unknown or blocked?\\
Evidence & Retained sources, histories, reproducible reports & Can a later reviewer reconstruct the decision without relying on memory?\\
\bottomrule
\end{longtable}

\section{Recording a method-comparison decision}
\label{sec:comparison-decision-record}

The main volume explains how a comparison affects a decision. This table retains
the fields needed to record the result consistently.

\begin{table}[htbp]
\centering\small
\caption{Decision record required after a substantive method comparison.}
\begin{tabularx}{\textwidth}{P{39mm}Y}
\toprule
Decision field & Required interpretation\\\midrule
Primary criterion & State the observed result and uncertainty against the frozen criterion.\\
Veto status & Identify severe errors, integrity failures and invalid comparisons before discussing speed.\\
Viable candidates & List methods that remain eligible for further evaluation; do not imply an unsupported ranking.\\
Main uncertainty & Distinguish reference ambiguity, sampling error, distribution shift and untested implementation.\\
Next justified action & Promote a bounded use, extend the study, repair a mechanism, or rebuild an invalid harness.\\
Limit of conclusion & State the source perimeter, task population and deployment claims not established.\\
\bottomrule
\end{tabularx}
\end{table}
